% GENERATED FILE. Edit canonical lessons, then run npm run generate:books.
% canonical-source-hash: fnv1a64:89494773a9cc8829
% canonical-lessons: KA-C413-latari, KA-C413-juju, KA-C413-pana, KA-C413-andajisu, KA-C413-sarasari

\chapter{Taking a Chance}
\label{ch:ka-taking-a-chance}
\hlchaptermodality{413}

\begin{chapteropening}
\textbf{By the end of this chapter:} \emph{I can say a lottery, gambling, a bet, to estimate and an average.}

Complete the last lesson of chapter 413: I can say a lottery, gambling, a bet, to estimate and an average.
\end{chapteropening}

\section[\texorpdfstring{\kn{ಲಾಟರಿ} (lāṭari)}{ಲಾಟರಿ (lāṭari)}]{\kn{ಲಾಟರಿ} (lāṭari) — a lottery}

\label{lesson:KA-C413-latari}



\begin{quote}
Before the new one: say the Kannada for hardship, then the Kannada for to ward off.
\end{quote}

\subsection*{You'll want to know: \kn{ಲಾಟರಿ}}
\textbf{\kn{ಲಾಟರಿ}} — \emph{lāṭari} — \textquotedblleft{}a lottery\textquotedblright{}.

\begin{quote}
\kn{ಅವನಿಗೆ} \kn{ಲಾಟರಿಯಲ್ಲಿ} \kn{ಹಣ} \kn{ಬಂತು}.
\end{quote}

\begin{quote}
\emph{avanige lāṭariyalli haṇa bantu.}
\end{quote}

\begin{quote}
He won money in the lottery.
\end{quote}

\kn{ಲಾಟರಿ} \kn{ಹೊಡೆ}, to win the lottery.

\begin{cousinweb}
\textbf{\kn{ಲಾಟರಿ}} (\emph{lāṭari}) is a loan from English \emph{lottery}.
\end{cousinweb}

\begin{grammarlens}[title={What you've built}]
One more word for saying what happens if.
\end{grammarlens}

\subsection*{Your turn}
\begin{itemize}
\raggedright
  \item \emph{Say it:} \emph{lāṭari}
  \item \emph{Say it:} \emph{lāṭari}, once more
  \item \emph{Say it:} \emph{avanige lāṭariyalli haṇa bantu.} — He won money in the lottery.
  \item \emph{Recall:} say the Kannada for hardship, then the Kannada for to ward off, then say \emph{lāṭari} again
\end{itemize}

\begin{tcolorbox}[breakable,colback=teal!4,colframe=teal!35!black,title={Before you move on}]
What does \kn{ಲಾಟರಿ} mean? (\textquotedblleft{}A lottery\textquotedblright{}.) Say it once more.
\end{tcolorbox}
\section[\texorpdfstring{\kn{ಜೂಜು} (jūju)}{ಜೂಜು (jūju)}]{\kn{ಜೂಜು} (jūju) — gambling}

\label{lesson:KA-C413-juju}



\begin{quote}
Before the new one: say the Kannada for to ward off, then the Kannada for a lottery.
\end{quote}

\subsection*{You'll want to know: \kn{ಜೂಜು}}
\textbf{\kn{ಜೂಜು}} — \emph{jūju} — \textquotedblleft{}gambling\textquotedblright{}.

\begin{quote}
\kn{ಜೂಜಿನಿಂದ} \kn{ಅವನು} \kn{ಮನೆ} \kn{ಕಳೆದುಕೊಂಡನು}.
\end{quote}

\begin{quote}
\emph{jūjininda avanu mane kaḷedukoṇḍanu.}
\end{quote}

\begin{quote}
He lost his house through gambling.
\end{quote}

\kn{ಜೂಜಾಡು}, to gamble.

\begin{cousinweb}
\textbf{\kn{ಜೂಜು}} (\emph{jūju}) is usually traced to Prakrit \emph{jūa}, \textquotedblleft{}gambling\textquotedblright{}, from Sanskrit \emph{dyūta}, \textquotedblleft{}gambling\textquotedblright{} (CDIAL 6608); Kannada is not among the forms CDIAL lists.
\end{cousinweb}

\begin{grammarlens}[title={What you've built}]
One more word for saying what happens if.
\end{grammarlens}

\subsection*{Your turn}
\begin{itemize}
\raggedright
  \item \emph{Say it:} \emph{jūju}
  \item \emph{Say it:} \emph{jūju}, once more
  \item \emph{Say it:} \emph{jūjininda avanu mane kaḷedukoṇḍanu.} — He lost his house through gambling.
  \item \emph{Recall:} say the Kannada for to ward off, then the Kannada for a lottery, then say \emph{jūju} again
\end{itemize}

\begin{tcolorbox}[breakable,colback=teal!4,colframe=teal!35!black,title={Before you move on}]
What does \kn{ಜೂಜು} mean? (\textquotedblleft{}Gambling\textquotedblright{}.) Say it once more.
\end{tcolorbox}
\section[\texorpdfstring{\kn{ಪಣ} (paṇa)}{ಪಣ (paṇa)}]{\kn{ಪಣ} (paṇa) — a bet; a stake}

\label{lesson:KA-C413-pana}



\begin{quote}
Before the new one: say the Kannada for a lottery, then the Kannada for gambling.
\end{quote}

\subsection*{You'll want to know: \kn{ಪಣ}}
\textbf{\kn{ಪಣ}} — \emph{paṇa} — \textquotedblleft{}a bet; a stake\textquotedblright{}.

\begin{quote}
\kn{ನಾಳೆ} \kn{ಮಳೆ} \kn{ಬರುತ್ತದೆ} \kn{ಅಂತ} \kn{ನಾವು} \kn{ಪಣ} \kn{ಕಟ್ಟಿದೆವು}.
\end{quote}

\begin{quote}
\emph{nāḷe maḷe baruttade anta nāvu paṇa kaṭṭidevu.}
\end{quote}

\begin{quote}
We bet that it would rain tomorrow.
\end{quote}

\kn{ಪಣ} \kn{ಕಟ್ಟು}, to place a bet; \kn{ಪಣಕ್ಕೆ} \kn{ಇಡು}, to put at stake.

\begin{cousinweb}
\textbf{\kn{ಪಣ}} (\emph{paṇa}) is borrowed from Sanskrit \emph{paṇa}, \textquotedblleft{}a wager, a stake\textquotedblright{} (CDIAL 7714).
\end{cousinweb}

\begin{grammarlens}[title={What you've built}]
One more word for saying what happens if.
\end{grammarlens}

\subsection*{Your turn}
\begin{itemize}
\raggedright
  \item \emph{Say it:} \emph{paṇa}
  \item \emph{Say it:} \emph{paṇa}, once more
  \item \emph{Say it:} \emph{nāḷe maḷe baruttade anta nāvu paṇa kaṭṭidevu.} — We bet that it would rain tomorrow.
  \item \emph{Recall:} say the Kannada for a lottery, then the Kannada for gambling, then say \emph{paṇa} again
\end{itemize}

\begin{tcolorbox}[breakable,colback=teal!4,colframe=teal!35!black,title={Before you move on}]
What does \kn{ಪಣ} mean? (\textquotedblleft{}A bet; a stake\textquotedblright{}.) Say it once more.
\end{tcolorbox}
\section[\texorpdfstring{\kn{ಅಂದಾಜಿಸು} (andājisu)}{ಅಂದಾಜಿಸು (andājisu)}]{\kn{ಅಂದಾಜಿಸು} (andājisu) — to estimate}

\label{lesson:KA-C413-andajisu}



\begin{quote}
Before the new one: say the Kannada for gambling, then the Kannada for a bet.
\end{quote}

\subsection*{You'll want to know: \kn{ಅಂದಾಜಿಸು}}
\textbf{\kn{ಅಂದಾಜಿಸು}} — \emph{andājisu} — \textquotedblleft{}to estimate\textquotedblright{}.

\begin{quote}
\kn{ಈ} \kn{ಕೆಲಸಕ್ಕೆ} \kn{ಹತ್ತು} \kn{ದಿನ} \kn{ಬೇಕು} \kn{ಅಂತ} \kn{ನಾನು} \kn{ಅಂದಾಜಿಸಿದೆ}.
\end{quote}

\begin{quote}
\emph{ī kelasakke hattu dina bēku anta nānu andājiside.}
\end{quote}

\begin{quote}
I estimated that this job needs ten days.
\end{quote}

Past (I …): \kn{ಅಂದಾಜಿಸಿದೆ}, \emph{andājiside}.

\kn{ಅಂದಾಜು}, an estimate; \kn{ಅಂದಾಜು} \kn{ಮಾಡು}, to guess.

\begin{cousinweb}
\textbf{\kn{ಅಂದಾಜಿಸು}} (\emph{andājisu}) is built on \emph{andāju}, \textquotedblleft{}an estimate\textquotedblright{}, a loan from Persian \emph{andāza}, \textquotedblleft{}a measure\textquotedblright{}, through Hindustani, with the Kannada verb ending -\emph{isu}.
\end{cousinweb}

\begin{grammarlens}[title={What you've built}]
One more word for saying what happens if.
\end{grammarlens}

\subsection*{Your turn}
\begin{itemize}
\raggedright
  \item \emph{Say it:} \emph{andājisu}
  \item \emph{Say it:} \emph{andājisu}, once more
  \item \emph{Say it:} \emph{ī kelasakke hattu dina bēku anta nānu andājiside.} — I estimated that this job needs ten days.
  \item \emph{Recall:} say the Kannada for gambling, then the Kannada for a bet, then say \emph{andājisu} again
\end{itemize}

\begin{tcolorbox}[breakable,colback=teal!4,colframe=teal!35!black,title={Before you move on}]
What does \kn{ಅಂದಾಜಿಸು} mean? (\textquotedblleft{}To estimate\textquotedblright{}.) Say it once more.
\end{tcolorbox}
\section[\texorpdfstring{\kn{ಸರಾಸರಿ} (sarāsari)}{ಸರಾಸರಿ (sarāsari)}]{\kn{ಸರಾಸರಿ} (sarāsari) — an average}

\label{lesson:KA-C413-sarasari}



\begin{quote}
Before the new one: say the Kannada for a bet, then the Kannada for to estimate.
\end{quote}

\subsection*{You'll want to know: \kn{ಸರಾಸರಿ}}
\textbf{\kn{ಸರಾಸರಿ}} — \emph{sarāsari} — \textquotedblleft{}an average\textquotedblright{}.

\begin{quote}
\kn{ತರಗತಿಯ} \kn{ಸರಾಸರಿ} \kn{ಅಂಕ} \kn{ಅರವತ್ತು}.
\end{quote}

\begin{quote}
\emph{taragatiya sarāsari aṅka aravattu.}
\end{quote}

\begin{quote}
The class's average mark is sixty.
\end{quote}

\kn{ಸರಾಸರಿ} \kn{ತೆಗೆ}, to work out the average.

\begin{cousinweb}
\textbf{\kn{ಸರಾಸರಿ}} (\emph{sarāsari}) is usually traced to Persian \emph{sarāsarī}, \textquotedblleft{}from end to end, general\textquotedblright{}, which reached Kannada through Hindustani or Marathi.
\end{cousinweb}

\begin{grammarlens}[title={What you've built}]
One more word for saying what happens if.
\end{grammarlens}

\subsection*{Your turn}
\begin{itemize}
\raggedright
  \item \emph{Say it:} \emph{sarāsari}
  \item \emph{Say it:} \emph{sarāsari}, once more
  \item \emph{Say it:} \emph{taragatiya sarāsari aṅka aravattu.} — The class's average mark is sixty.
  \item \emph{Recall:} say the Kannada for a bet, then the Kannada for to estimate, then say \emph{sarāsari} again
\end{itemize}

\begin{tcolorbox}[breakable,colback=teal!4,colframe=teal!35!black,title={Before you move on}]
What does \kn{ಸರಾಸರಿ} mean? (\textquotedblleft{}An average\textquotedblright{}.) Say it once more.
\end{tcolorbox}
